BFV, BGV, and CKKS parameter sets can require coefficient moduli larger than a machine word. The residue number system (RNS, \autoref{subsec:crt-application}) represents such coefficients by residues modulo smaller pairwise coprime moduli. This permits exact modular arithmetic using word-sized operations, without reducing the total modulus or changing the underlying security parameters. The required modulus sizes depend on the scheme, circuit, precision, and security target; no particular bit length is universal.

CRT makes modular addition and multiplication straightforward, but changing bases and implementing scale-and-round require additional algorithms. This section develops those primitives for BFV, BGV, and CKKS. Common TFHE implementations use torus coefficients represented within a machine word and therefore do not need this particular large-integer decomposition; that implementation choice is not a theorem forbidding RNS techniques in other TFHE variants.

In this section, $[x]_q=|x|_q$ denotes the unique integer representative in $[-q/2,q/2)$ (\autoref{subsec:modulo-centered}). Unless stated otherwise, RNS moduli are pairwise coprime odd integers; prime moduli are convenient but are not required by CRT. For odd $q$, the extreme representatives are $-(q-1)/2$ and $(q-1)/2$. Bounds on these integer representatives must be distinguished from congruences. This distinction is essential for both approximate and exact base conversion.

\begin{tcolorbox}[
    title = \textbf{Required Background},    % box title
    colback = white,    % light background; tweak to taste
    colframe = black,  % frame colour
    boxrule = 0.8pt,     % line thickness
    left = 1mm, right = 1mm, top = 1mm, bottom = 1mm % inner padding
]

\begin{itemize}
\item \autoref{sec:modulo}: \nameref{sec:modulo}
\item \autoref{sec:group}: \nameref{sec:group}
\item \autoref{sec:field}: \nameref{sec:field}
\item \autoref{sec:polynomial-ring}: \nameref{sec:polynomial-ring}
\item \autoref{sec:decomp}: \nameref{sec:decomp}
\item \autoref{sec:modulus-rescaling}: \nameref{sec:modulus-rescaling}
\item \autoref{sec:chinese-remainder}: \nameref{sec:chinese-remainder}
\item \autoref{sec:polynomial-interpolation}: \nameref{sec:polynomial-interpolation}
\item \autoref{sec:ntt}: \nameref{sec:ntt}
\item \autoref{sec:lattice}: \nameref{sec:lattice}
\item \autoref{sec:rlwe}: \nameref{sec:rlwe}
\item \autoref{sec:glwe}: \nameref{sec:glwe}
\item \autoref{sec:glwe-add-cipher}: \nameref{sec:glwe-add-cipher}
\item \autoref{sec:glwe-add-plain}: \nameref{sec:glwe-add-plain}
\item \autoref{sec:glwe-mult-plain}: \nameref{sec:glwe-mult-plain}
\item \autoref{subsec:modulus-switch-rlwe}: \nameref{subsec:modulus-switch-rlwe}
\item \autoref{sec:glwe-key-switching}: \nameref{sec:glwe-key-switching}
\item \autoref{sec:bfv}: \nameref{sec:bfv}
\item \autoref{sec:ckks}: \nameref{sec:ckks}
\item \autoref{sec:bgv}: \nameref{sec:bgv}
\end{itemize}
\end{tcolorbox}

\clearpage

\subsection{Fast Base Conversion: \textsf{FastBConv}}
\label{subsec:rns-fastbconv}

\noindent \textbf{- References:} \href{https://eprint.iacr.org/2016/510}{A Full RNS Variant of FV-like Somewhat Homomorphic Encryption Schemes}~\cite{rns-bfv}; \href{https://eprint.iacr.org/2022/657}{BASALISC: Programmable Hardware Accelerator for BGV FHE}~\cite{rns-bfv2}.

Let $q=\prod_{i=1}^kq_i$, $b=\prod_{j=1}^lb_j$, with all base moduli pairwise coprime. The input residues $x_i=[x]_{q_i}$ determine the centered integer $x=[x]_q$. Exact base conversion would output $([x]_{b_j})_{j=1}^l$. Fast base conversion omits one large reduction and instead uses the same unreduced CRT sum in every output coordinate:
\[
y_i=q/q_i,\qquad z_i=y_i^{-1}\bmod q_i,\qquad
T=\sum_{i=1}^k[x_iz_i]_{q_i}y_i.
\]
\begin{tcolorbox}[title={\textbf{\tboxlabel{\ref*{subsec:rns-fastbconv}} Fast Base Conversion: \textsf{FastBConv}}}]
\[
\textsf{FastBConv}(x,q,b)=([T]_{b_1},\ldots,[T]_{b_l}).
\]
There is an integer $u$ such that $T=x+uq$. For the centered odd-modulus convention above,
\[
|T|<kq/2,\quad |x|<q/2,\quad |u|<(k+1)/2,
\quad |u|\leq\lfloor k/2\rfloor.
\]
Thus the output represents $x+uq\pmod b$, not necessarily $x\pmod b$. Set $U_k=\lfloor k/2\rfloor$ for this bound.
\end{tcolorbox}

\begin{myproof}
Each summand has absolute value less than $q/2$, and the CRT congruence gives $T\equiv x\pmod q$. The bound on $u=(T-x)/q$ uses both $T$ and $x$; bounding $T$ alone does not bound $uq$ by the same interval. Every output residue is computed from this common $T$, so the error is a coherent multiple $uq$ across the destination base. Other residue conventions yield different bounds, including the unsigned convention in~\cite{rns-bfv}.
\end{myproof}

This is an approximation to base conversion, not a failure of CRT. Subsequent algorithms must remove the multiple of $q$, divide it into a controlled error, or explicitly include it in their correctness bound.

\subsection{Small Montgomery Reduction Algorithm: \textsf{SmallMont}}
\label{subsec:rns-smallmont}

Small Montgomery reduction limits the multiple of $q$ introduced by a fast base conversion. It does not generally make that conversion exact. The following centered-residue variant makes its assumptions explicit; the original unsigned-residue algorithm and its parameter bound appear in~\cite{rns-bfv}, Algorithm 2 and Lemma 4.

Let $x=[x]_q$, choose an odd auxiliary modulus $b_\alpha$ coprime to $q$ and $b$, and compute the fast CRT lift of $[b_\alpha x]_q$:
\[
T=[b_\alpha x]_q+uq=b_\alpha x+(v+u)q,
\qquad |u|\leq U_k.
\]
The input consists of the residues of this same $T$ modulo every $b_j$ and $b_\alpha$.
\begin{tcolorbox}[title={\textbf{\tboxlabel{\ref*{subsec:rns-smallmont}} Small Montgomery Reduction: \textsf{SmallMont}}}]
\begin{enumerate}
\item Compute $a=[Tq^{-1}]_{b_\alpha}$ using the supplied $b_\alpha$ residue.
\item For $j=1,\ldots,l$, output $r_j=[(T_{b_j}-qa)b_\alpha^{-1}]_{b_j}$.
\end{enumerate}
The result represents $r=(T-qa)/b_\alpha=x+u'q\pmod b$. If $b_\alpha>U_k+1/2$, then $u'\in\{-1,0,1\}$.
\end{tcolorbox}

\begin{myproof}
Because $a\equiv v+u\pmod{b_\alpha}$, the quotient $u'=(v+u-a)/b_\alpha$ is integral. Centered representatives give $|v|<(b_\alpha+1)/2$ and $|a|<b_\alpha/2$. Therefore
\[
|u'|<1+\frac{U_k+1/2}{b_\alpha}<2.
\]
This proves the stated three possible integer values. The numerator $T-qa$ is divisible by $b_\alpha$, and its quotient can be computed modulo each $b_j$ by multiplying by $b_\alpha^{-1}\bmod b_j$.
\end{myproof}

The auxiliary residue must remain available throughout the computation. Reconstructing $[T]_b$ and then reducing that value modulo $b_\alpha$ generally loses the original $b_\alpha$ residue. Also, $v$ can grow linearly with $b_\alpha$; a claim that it grows more slowly is not a proof of the bound.

\subsubsection{Improving \textsf{FastBConv} by Using \textsf{SmallMont}}
\label{subsubsec:rns-smallmont-fastbconv}

With the parameter condition in \autoref{subsec:rns-smallmont},
\[
\textsf{SmallMont}\bigl(\textsf{FastBConv}([b_\alpha x]_q,q,bb_\alpha),bb_\alpha,b_\alpha,q\bigr)
=[x+u'q]_b,\qquad u'\in\{-1,0,1\}.
\]
The preliminary multiplication by $b_\alpha$ is essential. The result is still an approximate lift of the centered $x$, and is distinct from \textsf{FastBConvEx} in \autoref{subsec:rns-fastbconvex}.


\subsection{RNS-based ModRaise: \textsf{ModRaise\textsubscript{RNS}}}
\label{subsec:rns-modraise}

\noindent \textbf{- Reference:} \href{https://eprint.iacr.org/2018/931}{A Full RNS Variant of Approximate Homomorphic Encryption}~\cite{rns-ckks}, Section 3.2.

Base extension from $q$ to $qb$ retains the residues modulo the $q_i$ and appends residues modulo the $b_j$. With fast base conversion it chooses an approximate integer lift:
\begin{tcolorbox}[title={\textbf{\tboxlabel{\ref*{subsec:rns-modraise}} \textsf{ModRaise\textsubscript{RNS}}}}]
\[
\textsf{ModRaise}_{\mathrm{RNS}}(x,q,qb)
=(x_1,\ldots,x_k,\textsf{FastBConv}(x,q,b)).
\]
The combined vector represents $x+uq\pmod{qb}$, with $|u|\leq U_k$. Using the Montgomery preprocessing of \autoref{subsec:rns-smallmont} reduces this to an appropriate lift $x+u'q$ with $|u'|\leq1$.
\end{tcolorbox}

This is a choice of integer representative, not a canonical ring embedding from $\mathbb Z_q$ into $\mathbb Z_{qb}$. If applied coefficient-wise to ciphertext components, write the new representatives as $A+qU_A$ and $B+qU_B$. Their phase changes by
\[
q(U_B+U_AS).
\]
Thus a component-wise bound $\|U_A\|_\infty,\|U_B\|_\infty\leq U_k$ gives the phase bound $qU_k(1+\|S\|_1)$, not simply $qU_k$. There can also be an existing multiple of $q$ in the integer phase $B+AS$. These terms must be handled by the consuming multiplication or bootstrapping algorithm; they are not independent fresh RLWE errors.

\subsection{RNS-based ModDrop: \textsf{ModDrop\textsubscript{RNS}}}
\label{subsec:rns-moddrop}

\textsf{ModDrop\textsubscript{RNS}} removes the residues for the factors of $b$ from a representation modulo $qb$:
\begin{tcolorbox}[title={\textbf{\tboxlabel{\ref*{subsec:rns-moddrop}} \textsf{ModDrop\textsubscript{RNS}}}}]
\[
(x_{q_1},\ldots,x_{q_k},x_{b_1},\ldots,x_{b_l})
\longmapsto(x_{q_1},\ldots,x_{q_k}).
\]
This is exact reduction modulo $q$; it performs no scaling or rounding.
\end{tcolorbox}

For CKKS or BGV, the same centered integer phase survives only if it already fits in $[-q/2,q/2)$ coefficient-wise. CKKS then keeps its scale and BGV keeps its plaintext modulus. Otherwise a wrapped phase can change the decoded message. For BFV, even an unchanged phase does not in general preserve the message, because the standard scale changes from $\lfloor qb/t\rfloor$ to $\lfloor q/t\rfloor$. For example, with initial modulus $105$, $t=7$, and noiseless phase $15$ encoding $m=1$, dropping to modulus $35$ leaves phase $15$ but standard BFV decoding gives $3$. Use a BFV-specific scale-and-round operation when message preservation is required.

\subsection{RNS-based Modulus Switch: \textsf{ModSwitch\textsubscript{RNS}}}
\label{subsec:rns-modswitch}

Let $x$ be an integer given in the combined bases of $q$ and $b$. For odd $b$, define $r=[x]_b$, so $(x-r)/b=\lfloor x/b\rceil$. The following operation approximately divides by $b$ and removes its base factors:
\begin{tcolorbox}[title={\textbf{\tboxlabel{\ref*{subsec:rns-modswitch}} \textsf{ModSwitch\textsubscript{RNS}}}}]
\begin{enumerate}
\item Convert the residues of $r$ from base $b$ to base $q$ using \textsf{FastBConv}. Its common integer lift is $r+ub$, with $|u|\leq U_l$.
\item For each $i=1,\ldots,k$, output
\[
y_i=[b^{-1}(x_{q_i}-(r+ub)_{q_i})]_{q_i}.
\]
\end{enumerate}
The result represents
\[
y=\left\lfloor\frac{x}{b}\right\rceil-u\pmod q,\qquad |u|\leq U_l.
\]
\end{tcolorbox}

This follows by substituting $x-r=b\lfloor x/b\rceil$. There is no need to assume $b\ll q$ for this congruence. For a single odd dropped modulus, $l=1$ and $u=0$, so the result is exact nearest-integer division. With several dropped factors, the additional error $-u$ must be included in the ciphertext phase bound. Applied to $(A,B)$, the total perturbation also includes multiplication of the component errors by $S$; it is not bounded solely by a single coefficient's $u$.

CKKS can use this as a rescaling primitive, tracking the actual scale division by $b$ and all rounding errors. BFV and BGV require their own plaintext-preservation rules. In particular, ordinary scale-and-round is not the BGV modulus switch of \autoref{subsec:rns-modswitch-bgv}.

\subsubsection{Comparing \textsf{ModSwitch\textsubscript{RNS}}, \textsf{ModRaise\textsubscript{RNS}}, and \textsf{ModDrop\textsubscript{RNS}}}

\textsf{ModRaise\textsubscript{RNS}} chooses a lift into a larger combined base; \textsf{ModDrop\textsubscript{RNS}} exactly reduces to a smaller sub-base; \textsf{ModSwitch\textsubscript{RNS}} additionally divides the represented integer approximately. These are different operations. Whether any of them preserves a ciphertext's plaintext depends on the scheme, its scale or plaintext modulus, and the full centered phase bound.




\subsection{RNS-based Decryption}
\label{subsec:rns-dec}

This subsection will explain how to efficiently decrypt RNS-based ciphertexts for BFV, CKKS, and BGV.


\subsubsection{BFV Decryption: $\textsf{Dec}_{\textsf{RNS}}^{\textsf{BFV}}$}
\label{subsubsec:rns-dec-bfv}

Use the book's ciphertext convention $B=-AS+\Delta M+E\pmod q$. The standard BFV decoding operation on each centered phase coefficient $x=[B+AS]_q$ is $[\lfloor tx/q\rceil]_t$; division by $\Delta=\lfloor q/t\rfloor$ is not algebraically identical. The following redundant-modulus correction implements scale-and-round when its rounding margin is satisfied, as in~\cite{rns-bfv}, Section 3.

Let $q=\prod_{i=1}^kq_i$ be odd, $\gcd(q,t)=1$, and choose an odd $\gamma$ coprime to $qt$. Put
\[
m_0=\left\lfloor\frac{tx}{q}\right\rceil,\quad
\delta=\frac{tx}{q}-m_0,\quad v=[t\gamma x]_q.
\]
The fast CRT lift of $v$ is $v+uq$, $|u|\leq U_k$. Compute its residues modulo $\gamma$ and $t$, and multiply by $-q^{-1}$ in each modulus to obtain $(y_\gamma,y_t)$. Then
\[
y\equiv-q^{-1}(v+uq)
\equiv\left\lfloor\frac{t\gamma x}{q}\right\rceil-u
=\gamma m_0+\eta\pmod{\gamma t},
\qquad \eta=\lfloor\gamma\delta\rceil-u.
\]
If $|\eta|<\gamma/2$, its centered residue is exactly $y_\gamma=\eta$, and therefore
\[
[m_0]_t=[(y_t-y_\gamma)\gamma^{-1}]_t.
\]
A sufficient condition is $\gamma|\delta|+1/2+U_k<\gamma/2$.

\begin{tcolorbox}[title={\textbf{\tboxlabel{\ref*{subsubsec:rns-dec-bfv}} $\textsf{Dec}_{\textsf{RNS}}^{\textsf{BFV}}$}}]
\begin{enumerate}
\item Evaluate the phase residues $x_i=[B+AS]_{q_i}$.
\item Choose $\gamma$ satisfying coprimality and the rounding-margin bound, and fast-convert $[t\gamma x]_q$ to the two moduli $\gamma,t$.
\item Multiply the two residues by $-q^{-1}$ in their respective moduli, obtaining $(y_\gamma,y_t)$ with centered $y_\gamma$.
\item Output $[(y_t-y_\gamma)\gamma^{-1}]_t$.
\end{enumerate}
\end{tcolorbox}

Correct BFV decryption additionally requires $[m_0]_t$ to equal the intended plaintext. Merely choosing any coprime $\gamma$ is insufficient, and the correction threshold is $\gamma/2$, not $\gamma$. Values arbitrarily close to a rounding boundary require a correspondingly stronger precision or correction strategy. This is not an unrestricted exact rounding primitive for arbitrary inputs and moduli. The alternative scale-and-round constructions in~\cite{rns-bfv-hps} also specify their precision and error conditions.


\subsubsection{CKKS and BGV Decryption}
\label{subsubsec:rns-dec-ckks-bgv}

Compute $[B+AS]_q$ in the current RNS base. CKKS decodes its centered integer representative with the tracked scale; BGV reduces that representative modulo the plaintext modulus. CRT reconstruction is one valid implementation. Dropping to a single small modulus before decoding is valid only when the entire centered phase fits within half of that smaller modulus. A ciphertext valid at its current level need not satisfy this stronger bound. When necessary, use scheme-correct rescaling or modulus switching with an error analysis, rather than assuming that a bare \textsf{ModDrop} to the lowest level preserves the message.

\subsection{BGV's RNS-based Modulus Switch: $\textsf{ModSwitch}_{\textsf{RNS}}^{\textsf{BGV}}$}
\label{subsec:rns-modswitch-bgv}

BGV modulus switching must preserve the plaintext residue, in addition to reducing the ciphertext modulus. The ordinary rounding primitive of \autoref{subsec:rns-modswitch} does not enforce this condition. Nor can one obtain an unrestricted rounding operation by calling $\textsf{Dec}_{\textsf{RNS}}^{\textsf{BFV}}$ with $t=\hat q$: if $\hat q$ divides $q_l$, the required inverse $q_l^{-1}\bmod\hat q$ does not exist.

Here is a direct RNS implementation for a common modulus-chain step. Let $q=Qw$, with $\gcd(w,Q)=\gcd(w,t)=1$ and $w\equiv1\pmod t$. For each centered coefficient $a$ of either ciphertext component, define
\[
d=t[t^{-1}a]_w,\qquad a'=\frac{a-d}{w}\pmod Q.
\]
Because $d\equiv a\pmod w$ and $d\equiv0\pmod t$, this is an integer division and $a'\equiv a\pmod t$ before reduction modulo $Q$. It is computed in each retained modulus $q_i\mid Q$ as
\[
a'_i=[w^{-1}(a_i-d)]_{q_i}.
\]
If $w$ has several RNS factors, the centered residue $[t^{-1}a]_w$ must be reconstructed or converted with an exact, bounded procedure; a single-factor step needs no large reconstruction.

Let the input phase be $M+tE$ and use the same overflow polynomial as in \autoref{subsec:bgv-modulus-switch}. The pre-reduction output phase is
\[
Y=\frac{M+tE}{w}-\frac{d_B+d_AS}{w}.
\]
Since $q=Qw\equiv Q\pmod t$, this phase is congruent to $M$ modulo $t$. A sufficient no-wrap condition is
\[
\frac{\|M+tE\|_\infty}{w}+\frac{t}{2}(1+\|S\|_1)<\frac Q2.
\]
Under this condition the output is a valid BGV ciphertext for the same plaintext. More general modulus chains can track a plaintext correction factor; the congruence assumption here is a stated special case, not a universal restriction on BGV.


\subsection{Exact Fast Base Conversion: \textsf{FastBConvEx}}
\label{subsec:rns-fastbconvex}

Unlike \textsf{SmallMont}, exact base conversion with a redundant modulus removes the entire CRT ambiguity when an explicit bound is satisfied. This is the Shenoy--Kumaresan style correction used in~\cite{rns-bfv}, Section 4.4, Lemma 6. Let $b=\prod_{j=1}^lb_j$ and suppose the residues of an integer $x$ are known modulo every $b_j$ and an additional coprime odd modulus $b_\alpha$. Write
\[
\widehat x=[x]_b=x+\mu b,\qquad
T=\widehat x+ub=x+(\mu+u)b,
\]
where $T$ is the common unreduced CRT sum used by both fast conversions and $|u|\leq U_l$. Require $|\mu+u|<b_\alpha/2$; for a known bound $|\mu|\leq\lambda$, the sufficient condition $b_\alpha>2(\lambda+U_l)$ can be used.

\begin{tcolorbox}[title={\textbf{\tboxlabel{\ref*{subsec:rns-fastbconvex}} Exact Fast Base Conversion: \textsf{FastBConvEx}}}]
\begin{enumerate}
\item Retain the base-$b$ residues of $x$, and separately its supplied residue $x_\alpha=[x]_{b_\alpha}$.
\item Compute $\gamma=[(\textsf{FastBConv}(\widehat x,b,b_\alpha)-x_\alpha)b^{-1}]_{b_\alpha}$.
\item Output $[\textsf{FastBConv}(\widehat x,b,q)-\gamma b]_q$, coordinate-wise in the target base.
\end{enumerate}
\end{tcolorbox}

\begin{myproof}
Both conversions use $T$. Hence $\gamma=[\mu+u]_{b_\alpha}=\mu+u$ under the stated bound, and $T-\gamma b=x$. The output is exactly $[x]_q$. This is an algebraic identity; no fresh approximation error remains. If the bound is violated, the correction can instead differ by a multiple of $bb_\alpha$, which need not vanish modulo $q$.
\end{myproof}

\para{Necessity of the Centered (i.e., Signed) Residue Representation:} The correction $\gamma$ must be interpreted as the signed bounded integer $\mu+u$. The centered convention supplies this interpretation when $|\mu+u|<b_\alpha/2$. An unsigned implementation is also possible if it explicitly recovers the sign; simply treating an unsigned residue as that integer is incorrect. The requirement is a quantified range bound, not the unquantified statement that $x$ must be ``much smaller'' than its modulus.


\subsection{Decomposing Multiplication: \textsf{DecompMult\textsubscript{RNS}}}
\label{subsec:rns-decompmult}

RNS supplies an integer gadget for key switching and relinearization. With $q=\prod_iq_i$, define
\[
g_i=\frac q{q_i}\left[\left(\frac q{q_i}\right)^{-1}\right]_{q_i},\qquad
A_i=[A]_{q_i},\qquad A\equiv\sum_iA_ig_i\pmod q.
\]
Here each $A_i$ is lifted coefficient-wise to a small integer polynomial. When multiplying it by a ciphertext in $R_q$, reduce those same small integers in \emph{every} modulus of the RNS base. Keeping only the $i$-th coordinate and filling the others with zero would give a different multiplier and a different error bound.

\begin{tcolorbox}[title={\textbf{\tboxlabel{\ref*{subsec:rns-decompmult}} \textsf{DecompMult\textsubscript{RNS}}}}]
For switching $S$ to $S'$, generate evaluation keys $(a_i,b_i)$ satisfying
\[
b_i+a_iS'\equiv g_iS+E_i\pmod q.
\]
Then $\sum_iA_i(a_i,b_i)$ has phase $AS+\sum_iA_iE_i$ modulo $q$. Adding $(0,B)$ completes key switching. In BGV use $tE_i$ instead of $E_i$ so that the added error is divisible by the plaintext modulus.

For relinearization, replace $g_iS$ by $g_iS^2$ in the keys, decompose $D_2$, and add the result to $(D_1,D_0)$ in the book's $(A,B)$ ordering. The resulting phase is
\[
D_0+D_1S+D_2S^2+\sum_iD_{2,i}E_i\pmod q.
\]
\end{tcolorbox}

For negacyclic polynomials, $\|A_iE_i\|_\infty\leq n\|A_i\|_\infty\|E_i\|_\infty$, so an explicit bound on the added error is $n\sum_i(q_i/2)\|E_i\|_\infty$. One can instead put the modular inverse into the digit $[A_i(q/q_i)^{-1}]_{q_i}$ and use weight $q/q_i$, as in~\cite{rns-bfv}, Section 4.5. Both choices are valid, but their digits and keys must not be mixed. Encryption symbols in these identities denote phase relations, not equalities of freshly sampled ciphertext distributions.


\subsection{Applying RNS Techniques to FHE Operations}
\label{subsec:rns-application}

This subsection will explain how the RNS primitives we have learned so far are used to handle FHE operations for RNS-based ciphertexts in BFV, CKKS, and BGV.


\subsubsection{Addition and Multiplication of Polynomials}
\label{subsubsec:rns-application-basic}

RNS decomposes the coefficient modulus, so scalar addition and scalar multiplication are performed separately modulo each $q_i$ (\autoref{subsec:crt-application}). Polynomial addition is coefficient-wise, but polynomial multiplication still requires negacyclic convolution modulo $X^n+1$. It becomes pointwise multiplication only after an appropriate NTT, when the modulus supports that transform. The following scalar identities describe the coefficient arithmetic used inside those polynomial operations.

$P^{\langle 1 \rangle} = \sum\limits_{a=0}^{n-1} c^{\langle 1 \rangle}_a \cdot X^a \in \mathcal{R}_{\langle n, q\rangle}$

$P^{\langle 2 \rangle} = \sum\limits_{b=0}^{n-1} c^{\langle 2 \rangle}_b \cdot X^b \in \mathcal{R}_{\langle n, q\rangle}$

$ $

In the RNS-variant FHE schemes, we express each coefficient of the polynomial as an RNS residue vector as follows:

$c^{\langle 1 \rangle}_a = (c^{\langle 1 \rangle}_{a,1}, c^{\langle 1 \rangle}_{a,2}, \cdots, c^{\langle 1 \rangle}_{a,k}) \in \mathbb{Z}_{q_1} \times \mathbb{Z}_{q_2} \times \cdots \times \mathbb{Z}_{q_k}$ \textcolor{red}{ $\rhd$ for $a \in [0, n - 1]$}

$c^{\langle 2 \rangle}_b = (c^{\langle 2 \rangle}_{b,1}, c^{\langle 2 \rangle}_{b,2}, \cdots, c^{\langle 2 \rangle}_{b,k}) \in \mathbb{Z}_{q_1} \times \mathbb{Z}_{q_2} \times \cdots \times \mathbb{Z}_{q_k}$ \textcolor{red}{ $\rhd$ for $b \in [0, n - 1]$}

$ $

Given the above RNS setup, when we add or multiply two polynomials, each coefficient-to-coefficient addition is computed as element-wise additions of two RNS residue vectors as follows:

$c^{\langle 1 \rangle}_a + c^{\langle 2 \rangle}_b \equiv \sum\limits_{i=1}^k (c^{\langle 1 \rangle}_{a,i} + c^{\langle 2 \rangle}_{b,i}) y_i z_i \bmod q $ \textcolor{red}{ $\rhd$ where $y_i = \dfrac{q}{q_i}$, $z_i = |y_i^{-1}|_{q_i}$}

$\Longleftrightarrow (c^{\langle 1 \rangle}_{a,1} + c^{\langle 2 \rangle}_{b,1}, \text{ } c^{\langle 1 \rangle}_{a,2} + c^{\langle 2 \rangle}_{b,2}, \cdots, \text{ } c^{\langle 1 \rangle}_{a,k} + c^{\langle 2 \rangle}_{b,k}) \in \mathbb{Z}_{q_1} \times \mathbb{Z}_{q_2} \times \cdots \times \mathbb{Z}_{q_k}$

$ $

Similarly, each coefficient-to-coefficient multiplication is computed as element-wise multiplications of two RNS residue vectors as follows:

$c^{\langle 1 \rangle}_a \cdot c^{\langle 2 \rangle}_b \equiv \sum\limits_{i=1}^k (c^{\langle 1 \rangle}_{a,i} \cdot c^{\langle 2 \rangle}_{b,i}) y_i z_i \bmod q $

$\Longleftrightarrow (c^{\langle 1 \rangle}_{a,1} \cdot c^{\langle 2 \rangle}_{b,1}, \text{ } c^{\langle 1 \rangle}_{a,2} \cdot c^{\langle 2 \rangle}_{b,2}, \cdots, \text{ } c^{\langle 1 \rangle}_{a,k} \cdot c^{\langle 2 \rangle}_{b,k}) \in \mathbb{Z}_{q_1} \times \mathbb{Z}_{q_2} \times \cdots \times \mathbb{Z}_{q_k}$

$ $

Using the above isomorphism, we can efficiently compute ciphertext-to-plaintext addition, ciphertext-to-ciphertext addition, and ciphertext-to-plaintext multiplication of big polynomial coefficients (e.g., 1000 bits big) based on small RNS residues (e.g., 30 bits each).


\subsubsection{Key Switching}
\label{subsubsec:rns-application-key-switching}

With the evaluation keys and lifted digits of \autoref{subsec:rns-decompmult}, compute
\[
(A',B')=(0,B)+\textsf{DecompMult}_{\mathrm{RNS}}(A,S_{\langle S',\mathrm{RNS}\rangle}).
\]
If the input phase is $B+AS$, the output phase is $B+AS+E_{\mathrm{ks}}$ modulo $q$, with $E_{\mathrm{ks}}=\sum_iA_iE_i$. BFV and CKKS retain their current message scale while adding this error. BGV uses evaluation keys with errors divisible by $t$, so its phase remains congruent to the same message modulo $t$. All operations can be performed in the current RNS base, but the phase bounds remain scheme-specific.


\subsubsection{Input Slot Rotation}
\label{subsubsec:rns-rotation}

For $J(h)=5^h\bmod2n$, apply the automorphism $X\mapsto X^{J(h)}$ to both ciphertext components. Reduce exponents modulo $2n$, with $X^{n+i}=-X^i$. In coefficient representation this permutes entire RNS residue vectors and negates the relevant vectors in every modulus. It introduces no approximation error by itself.

The resulting phase is the original phase evaluated at $X^{J(h)}$, under secret $S(X^{J(h)})$. Then use the RNS key-switching operation in \autoref{subsubsec:rns-application-key-switching} to return to $S(X)$. This step adds key-switching error, with a factor of $t$ for BGV. The induced slot permutation depends on the chosen root ordering: CKKS has $n/2$ independent complex slots, while fully split BFV/BGV batching has two orbits of $n/2$ scalar slots under powers of $5$. RNS does not change that permutation.

\subsubsection{BFV's Ciphertext-to-Ciphertext Multiplication}
\label{subsubsec:rns-cipher-mult-bfv}

The BEHZ RNS multiplication in~\cite{rns-bfv}, Sections 4.2--4.5, performs base extension, tensor multiplication, approximate scale-and-round, exact conversion back, and relinearization. Its auxiliary-base and error bounds are part of the algorithm. The HPS construction~\cite{rns-bfv-hps}, Sections 2.2--2.4 and 4.5, uses different conversion/scaling techniques; the variants should not be conflated.

The following states the centered-residue version of the primitives used here. Let $b=\prod_jb_j$. Use a redundant modulus $b_\alpha$ for the final exact conversion and a \emph{distinct} temporary Montgomery modulus $\widetilde m$ for controlling the initial lifts. All base moduli must satisfy the required coprimality conditions.
\begin{enumerate}
\item \textbf{Base extension.} For each component coefficient $x$, fast-convert $[\widetilde m x]_q$ to the combined auxiliary base $bb_\alpha\widetilde m$, then apply \textsf{SmallMont} with temporary modulus $\widetilde m$. This leaves residues modulo $bb_\alpha$ for a lift $\widehat x=x+uq$, $|u|\leq1$, and retains the original base-$q$ residues. The auxiliary residue $b_\alpha$ is kept for the later exact conversion; it is not discarded as the Montgomery modulus.
\item \textbf{Tensor multiplication.} In the combined base $qbb_\alpha$, compute the negacyclic integer-polynomial products
\[
D_0=\widehat B_1\widehat B_2,\quad
D_1=\widehat A_1\widehat B_2+\widehat A_2\widehat B_1,\quad
D_2=\widehat A_1\widehat A_2.
\]
Choose the auxiliary range so that the intended integers $tD_j$ can be represented without wraparound, for example by verifying $\|tD_j\|_\infty<qbb_\alpha/2$. The coefficient convolution, $n$, $t$, and the component-lift bounds all matter; $b>\Delta$ alone is insufficient.
\item \textbf{Approximate scale-and-round.} Multiply each $D_j$ by $t$ and apply \textsf{ModSwitch\textsubscript{RNS}} dropping the base of $q$. The output represents
\[
Z_j=\left\lfloor\frac{tD_j}{q}\right\rceil-u_j
=\frac{tD_j}{q}+\rho_j,\quad
\|u_j\|_\infty\leq U_k,\quad
\|\rho_j\|_\infty\leq U_k+\frac12.
\]
in base $bb_\alpha$. This is approximate scale-and-round, not exact rounding.
\item \textbf{Exact return to base $q$.} Apply \textsf{FastBConvEx} to each $Z_j$ using its retained $b_\alpha$ residue. For every coefficient, the bound $|\mu+u|<b_\alpha/2$ in \autoref{subsec:rns-fastbconvex} must hold. An auxiliary-range bound for $tD_j$ does not replace this separate correction bound.
\item \textbf{Relinearization.} Compute $(Z_1,Z_0)+\textsf{DecompMult}_{\mathrm{RNS}}(Z_2,S^2_{\langle S,\mathrm{RNS}\rangle})$ using the $(A,B)$ ordering. This adds the evaluation-key error $E_{\mathrm{ks}}$.
\end{enumerate}

Let $H_i=\widehat B_i+\widehat A_iS$ denote the lifted integer phases. The output phase satisfies
\[
F\equiv\frac tqH_1H_2+\rho_0+\rho_1S+\rho_2S^2+E_{\mathrm{ks}}\pmod q.
\]
One cannot replace the three component roundings by a single rounding of $H_1H_2$. The bound must include the secret-dependent terms, input errors, lift multiples, and the discrepancy between $t\Delta^2/q$ and $\Delta$. For example, writing a centered output phase as $\Delta M_{\mathrm{out}}+E_{\mathrm{out}}$, a sufficient decoding condition is
\[
\|\Delta M_{\mathrm{out}}+E_{\mathrm{out}}\|_\infty<q/2,\qquad
\left\|\frac tqE_{\mathrm{out}}+\left(\frac{t\Delta}{q}-1\right)M_{\mathrm{out}}\right\|_\infty<1/2.
\]
Here $M_{\mathrm{out}}=[M_1M_2]_t$. These are absolute error conditions; comparing the error only with the message magnitude is invalid, especially when a message coefficient is zero. See \autoref{subsubsec:bfv-mult-cipher-summary} and the cited papers for the corresponding full parameter analysis.


\subsubsection{CKKS's Ciphertext-to-Ciphertext Multiplication}
\label{subsubsec:rns-cipher-mult-ckks}

Compute the tensor product directly in the current base, relinearize using the appropriate RNS evaluation key, and rescale when required by the scale/noise budget (\autoref{subsubsec:ckks-mult-cipher-summary}). The product phase must fit the current modulus before an approximate interpretation is used; a level number alone does not establish this.

If the input scales are $\Delta_1,\Delta_2$ and the dropped modulus is $w$, the new scale is $\Delta_1\Delta_2/w$. A single-prime rescale implements component-wise nearest-integer division exactly, but still contributes phase error $R_B+R_AS$. A multi-factor approximate conversion can contribute additional bounded errors. The BFV-specific factor $t$, wide tensor-product base, and final exact return to a fixed BFV modulus are not part of this CKKS procedure. Hybrid key-switching may nevertheless use its own auxiliary base.


\subsubsection{BGV's Ciphertext-to-Ciphertext Multiplication}
\label{subsubsec:rns-cipher-mult-bgv}

Tensor multiplication and RNS relinearization give the same plaintext product modulo $t$, with an error still divisible by $t$ (\autoref{subsec:bgv-mult-cipher}). Evaluation-key errors must also be multiples of $t$. If modulus switching is needed, use $\textsf{ModSwitch}_{\mathrm{RNS}}^{\mathrm{BGV}}$ from \autoref{subsec:rns-modswitch-bgv}, or a more general congruence-preserving variant with an explicitly tracked correction factor. Ordinary CKKS rescaling is not a substitute. The factor $t$ remains fixed; correctness and useful noise reduction depend on the full output-phase bound.


\subsubsection{BFV's Bootstrapping}
\label{subsubsec:rns-bfv-bootstrapping}

RNS is a representation of the ciphertext coefficient modulus; it does not prohibit a plaintext modulus $p^\varepsilon$. A prime power can also be one CRT modulus, provided it is coprime to the other base moduli. It is only invalid to treat repeated factors $p,p,\ldots,p$ as distinct coprime CRT coordinates.

The BFV bootstrapping construction of~\cite{cryptoeprint:2022/1363}, Sections 3.2--3.4, uses the public preprocessing
\[
c_i'=\left[\left\lfloor\frac{p^\varepsilon}{q}c_i\right\rceil\right]_{p^\varepsilon},
\]
then homomorphically computes $[c_0'+c_1'S]_{p^\varepsilon}$ with the bootstrapping key. Here $(c_0,c_1)=(B,A)$, and the plaintext modulus for that homomorphic computation is $p^\varepsilon$. Linear transforms and digit extraction implement the rounded decryption map to plaintext modulus $p^r$, followed by the inverse transform. See \autoref{subsubsec:bfv-bootstrapping-summary} for the scheme-level discussion.

This coefficient-to-slot construction requires $\gcd(p,\mu)=1$ for the cyclotomic index $\mu$ (Section 2.3 of the cited paper). For $\mu=2n$ with power-of-two $n$, use odd $p$. A binary digit-extraction rule cannot by itself supply the missing plaintext-ring CRT in that ring.

Public preprocessing can reconstruct ciphertext coefficients from RNS and round exactly, or use a proven RNS scale-and-round routine with sufficient precision. Ciphertext arithmetic in the homomorphic circuit can then use a separate RNS base for its working ciphertext modulus. If approximate base conversion is used, its error must be included in the rounded-decryption bound before claiming correctness. The BEHZ/HPS multiplication routines alone do not prove a new bootstrapping construction, and the invalid argument about non-coprime factors does not justify replacing this preprocessing by an unproved chain of base extensions and reinterpretations.

For the idealized input phase $(q/p^r)M+E$, component rounding contributes $R_B+R_AS$, and the rounded-decryption error includes $(p^\varepsilon/q)E+R_B+R_AS$. If the actual encryption scale is $\lfloor q/p^r\rfloor$, its discrepancy must be included in $E$ in this expression. Recovering $M$ requires an explicit bound below $p^{\varepsilon-r}/2$, as well as valid intermediate ciphertexts throughout the homomorphic circuit.

\subsubsection{CKKS's Bootstrapping}
\label{subsubsec:rns-ckks-bootstrapping}

Use the CKKS sequence in \autoref{subsubsec:ckks-bootstrapping-summary}: modulus raising, coefficient-to-slot transformation, approximate modular reduction, and slot-to-coefficient transformation. Sparse bootstrapping additionally needs the partial-sum operation explained in \autoref{subsubsec:ckks-bootstrapping-time-reduction}. There is no separate explicit homomorphic secret-key inner product in this basic CKKS construction, nor a universal final reinterpretation step.

If approximate RNS lifting changes the components by $q_0U_A,q_0U_B$, the lifted phase is
\[
\Delta M+E+q_0(K+U_B+U_AS).
\]
The additional coefficient-wise lift bound is amplified by the secret polynomial. Even if each coefficient of $U_A,U_B$ lies in $\{-1,0,1\}$, the corresponding phase multiple need not. The approximation interval for the periodic modular-reduction circuit must cover this total overflow, and its error budget must include the linear transforms, rescaling, key switching, polynomial approximation, and retained input error. Periodicity alone does not make a finite homomorphic polynomial an exact modular-reduction function.


\subsubsection{BGV's Bootstrapping}
\label{subsubsec:rns-bgv-bootstrapping}

Follow the BGV construction in \autoref{subsec:bgv-bootstrapping}. First perform a valid modulus switch to $\hat q\equiv1\pmod{p^\varepsilon}$ and compute centered public coefficients $A'=[p^{\varepsilon-r}A]_{\hat q}$, $B'=[p^{\varepsilon-r}B]_{\hat q}$. The special modulus need not be a member of the current RNS base: exact public reconstruction and congruence-preserving switching are a valid implementation route.

Homomorphically evaluate $B'+A'S$ using a bootstrapping key at plaintext modulus $p^\varepsilon$ and a sufficiently large working ciphertext modulus $Q$. RNS represents $Q$; it does not require the plaintext modulus to equal an RNS factor. Then implement the appropriate coefficient/slot transforms and balanced rounded digit extraction, with a verified bound on $K'$ and on every homomorphic intermediate phase.

During digit extraction, each justified exact division by $p$ changes the plaintext modulus $p^j\rightarrow p^{j-1}$; after $\varepsilon-r$ such divisions it is $p^r$. Do not append another multiplication by $p^{-(\varepsilon-r)}$, which would divide the message twice. Also, the base-extension primitive $q\rightarrow qb$ does not directly describe an arbitrary conversion $\hat q\rightarrow Q$ when $\hat q$ does not divide $Q$. An alternative bootstrapping construction using such lifts needs its own correctness proof; it is not established by the RNS representation alone.


\subsubsection{Noise Impact of RNS Operations}
\label{subsubsec:rns-noise}

Track the actual effect of each primitive. \textsf{FastBConv} chooses a lift differing by a bounded multiple of the source modulus; \textsf{SmallMont} limits that multiple but does not generally eliminate it. \textsf{ModSwitch\textsubscript{RNS}} turns a base-conversion discrepancy into a bounded additive error after division. \textsf{FastBConvEx} is exact under its range bound and does not need \textsf{SmallMont} to remove an error. \textsf{ModDrop\textsubscript{RNS}} is exact residue reduction, but can invalidate a ciphertext if the new phase or scale conditions fail.

For component perturbations $\delta_A,\delta_B$, the induced phase perturbation is $\delta_B+\delta_AS$, with coefficient bound $\|\delta_B\|_\infty+\|\delta_A\|_\infty\|S\|_1$. Higher-degree ciphertexts similarly involve powers of $S$. Whether these perturbations are ultimately canceled, approximately removed, or retained as decryption error depends on the consuming scheme and operation; no single blanket noise statement applies to all RNS primitives.

\subsubsection{Python Source Code of RNS Primitives}
\label{subsubsec:rns-source-code}

We provide a \href{https://github.com/wang2000star/fhe-textbook-cn/blob/chinese-edition/source%20code/rns.py}{\underline{Python script}} implementing the following exemplary RNS primitives:

\noindent\textsf{FastBConv},
\textsf{ModRaise\textsubscript{RNS}},
\textsf{ModDrop\textsubscript{RNS}},
\textsf{ModSwitch\textsubscript{RNS}},
\textsf{SmallMont}, and
\textsf{FastBConvEx}.





